\subsubsection{The integral restriction at the terminal classes}
\label{subsec:terminal-realization}

The last restriction in the chain is a branching map between interlacing types. Its matrix must be evaluated in the actual class frames before composing the finite kernels.

Fix a word and either the reference tuple or one of its used
twists. Write \(N=n+1\) for the last two ranks. At embedding
\(\tau\) the two true weights \(\lambda,\sigma\) interlace,
and \(\lambda_N=\sigma_n\). Indeed the row weight is the
interlacing row of the last Harish--Chandra parameter, and
the shift in both ranks is \(-m(t+1)/2\), with \(t\) odd.
We use the polynomial weights after subtracting the common
last part \(q\). The actions on full weights include the
characters \(\det^q\).

\begin{proposition}[Integral terminal branching]\label{prop:integral-terminal-branching}
For interlacing weights with the same last part, the
normalized polynomial lowering operator below intertwines
the branching spaces and carries the integral Weyl
lattice into the indicated integral row lattice on
the standard frames, simultaneously at all embeddings.
Writing \(q\) for the common last part, the same determinant twist restores the two full weights
\[
 q=\lambda_{n+1}=\sigma_n,\qquad \chi(g)=\det(g)^q.
\]

\end{proposition}

\begin{proof}
In a single matrix of variables, let \(\nabla_\eta\) be the
integral span of products of flag minors of diagram
\(\eta-\eta_{\rm last}\), and \(h_\eta\) the product with all
row sets initial. Use the Weyl lattice \(\Delta_\eta\) generated
from \(h_\eta\) by divided powers of row operators.

On polynomial
models in opposite frames their perfect pairing is
\[
 (f,g)_\eta=
 \frac{\sum_\nu \nu!\,[X^\nu]f\,[X^\nu]g}
      {\sum_\nu \nu!\,([X^\nu]h_\eta)^2}.
\]
This is integral Weyl--dual Weyl duality for the split general
linear group over \(\mathbb Z\), applied to the Schur module
of the indicated partition. The normalization follows by pairing
the two highest polynomial vectors to one. Equivalently the
semistandard bideterminants are an integral basis for the flag
span. Pairing them with the divided row words for the same
tableaux gives a unitriangular matrix in tableau order. This
is the standard monomial proof of duality and shows why extending
the coefficient ring preserves it. On several embeddings take
the tensor product of these pairings. \label{cor:terminal-realization-1}

Put \(a_i=\sigma_i-\lambda_{i+1}\), \(b_i=\lambda_i-\sigma_i\)
for \(1\le i\le n\). Let \(D_i\) be the initial \(i\)-minor,
and \(D_i^*\) the minor in the first \(i\) columns with rows
\(1,\ldots,i-1,N\). The vector
\[
              \xi=\prod_{i=1}^n D_i^{a_i}(D_i^*)^{b_i}
\]
is nonzero and is a highest vector of small weight \(\sigma\)
in the restriction of the large representation. This follows
from alternation in the first \(n\) rows and the row counts,
including the determinant character. The classical interlacing
branching rule in characteristic zero makes this summand simple
of multiplicity one. Normalize its embedding \(J_0\) by
\(J_0 h_\sigma=\xi\).

Here the integral map in the period has a specified scalar.
Define \(d_0=\sum_i a_i\) and
\[
 R_i=[\lambda_{i+1},\min(\lambda_i,\lambda_{i+1}+d_0)]\cap\mathbb Z,
 \qquad
 z(\nu)=\prod_{i=1}^n\prod_{v\in R_i\setminus\{\sigma_i\}}
                      \prod_{j=1}^n(v-i-\nu_j+j).
\]
Let \(c=z(\sigma)\). The shifted interlacing intervals
\([\lambda_{i+1}-i,\lambda_i-i]\) are pairwise disjoint.
Thus \(c\ne0\). If \(\nu\) interlaces and
\(\sum\nu_i=\sum\sigma_i\), then every \(\nu_i\) belongs to
\(R_i\). Unless \(\nu=\sigma\), substituting \(v=\nu_i\) for
a differing coordinate gives \(z(\nu)=0\). The Capelli
characteristic polynomial for the small row action gives a
central operator with eigenvalue 
\[
  \prod_j(x-\nu_j+j)
\]
on highest weight \(\nu\).

In the usual column determinant
convention evaluate the Capelli determinant of
\(E+\operatorname{diag}(n-1,\ldots,0)-uI\) at \(u=x+n\)
and multiply by \((-1)^n\). The eigenvalue is the one just
stated. \cite{molev-branching} The row Lie operators \(E_{ij}\), with the integral
determinant twist on the diagonal, preserve the integral flag
span by differentiation of the integral action. Taking the
product at \(x=v-i\) gives an integral operator \(Z\) acting
by \(z(\nu)\). On the subspace of last-coordinate weight
\(\sum\lambda_i-\sum\sigma_i\), this means
\[
       Z|_{V_\nu}=z(\nu)\,\mathrm{id}_{V_\nu}
               =c\,\mathbf1_{\nu=\sigma}\,\mathrm{id}_{V_\nu}.
\]
The multiplicity-one summand \(V_\sigma\) is generated by
\(\xi\). Thus \(Z\) is \(c\) times that projection.

We verify that \(J=cJ_0\) takes \(\nabla_\sigma\) to
\(\nabla_\lambda\). Label the factors in a product \(P\) of
small flag minors. For each \(i\ge2\) append row \(N\) to
\(b_i\) columns of height \(i-1\) in the tableau for \(P\).
There are at least that many since

\[
  \sigma_{i-1}-\sigma_i\ge\lambda_i-\sigma_i.
\]
Use \(b_1\) additional singleton columns with entry \(N\).
For each enlarged minor use the first \(i\) matrix columns,
and keep the other minors. This gives an integral polynomial
\(M(P)\) of large shape and the required last weight. The
choice of which labeled columns to enlarge is fixed as the
row indices vary. The initial product maps to \(\xi\).
Cauchy--Binet shows that applying any small elementary row
substitution to \(P\) and \(M(P)\) gives expansions with the
same coefficients indexed by the same changes of the small
row tuples. The appended row is fixed. This remains true for
each coefficient in the substitution parameter and hence for
all sequences in divided row operators.

For each tuple let \(D(P)\) be the error term in the \(\sigma\)-summand. The substitution preserves the small diagonal torus weight, so whenever the difference is nonzero it satisfies
\[
  D(P)=ZM(P)-cJ_0P,\qquad \operatorname{wt}(D(P))=\operatorname{wt}(P).
\]
If this weight differs from the highest weight, its Fock pairing
with \(\xi\) vanishes by torus invariance of the Hermitian
version of the pairing.

At highest weight the row counts force
every nonzero small column to use rows \(1,\ldots,k\) for its
height \(k\). Count row 1 first, then row 2 in all columns
of height at least two, and continue. The signs of their
permutations agree after appending \(N\); a repeated index
kills both products. Hence at this weight the difference
vanishes. Put \(D(P)=ZM(P)-cJ_0P\). The simultaneous expansions give,
for every word \(T\) in the divided row operators,
\[
              0=\langle T D(P),\xi\rangle
               =\langle D(P),T^*\xi\rangle .
\]
Adjunction interchanges the two row indices. The vectors
\(T^*\xi\) span the summand containing \(D(P)\), so \(D(P)=0\). This proves
integrality without requiring the unprojected lift \(M\) to
respect the straightening relations.

For instance at \(\lambda=(2,1,0)\), \(\sigma=(2,0)\) one
has \(c=4\). Write \(x_i=X_{i1}\) and \(D_{ij}\) for the
two-column minor on rows \(i,j\). Then
\(J_0(x_1^2)=x_1D_{13}\), and applying the row derivation
from 1 to 2 gives
\[
 J_0(x_1x_2)=(x_2D_{13}+x_1D_{23})/2 .
\]
Its numerator has odd coefficients. The unscaled embedding
already fails to preserve the flag span at two in this example.

Take the normalized transpose of \(J\). It maps
\(\Delta_\lambda\) into \(\Delta_\sigma\) by the perfect
pairing. For the coefficient extraction, set the last column to zero and take the operator
\[
  P\longmapsto\Big[\prod_i X_{Ni}^{b_i}\Big]\big(P\big|_{\text{last column}=0}\big).
\]
 On \(\xi\) the result is
\(h_\sigma\). Starting with the largest \(i\), the required
power of \(X_{Ni}\) uses all the diagonal last-row terms of
the remaining minors of height \(i\). Project the resulting
polynomial orthogonally to the small representation in the
Fock form, and multiply by \(c(\xi,\xi)_\lambda\).
This is the formula for \(\mathrm{br}_{\rm std}\). Coefficient
extraction commutes with the small row action; the projection
commutes by unitary invariance and complexification. Both
intertwiners agree on \(\xi\), where the value is
\(c(\xi,\xi)_\lambda h_\sigma\), so multiplicity one makes
them equal. All coefficients descend to the rational
polynomial model. Tensoring proves integrality in the
standard frames simultaneously at all embeddings.
\label{cor:terminal-realization-2}
\end{proof}

\paragraph{Lowering in the class frame.}

Restriction on highest-weight spaces is rational. Integrality of its values must be checked in the frames of the two matched terminal classes.

\begin{lemma}[Integral terminal values in the class frame]
\label{lem:terminal-integral-class-frame}
At a lifted class and its matched terminal class, write \(L_{\rm st}\) for the integral coefficient lattice in the stage frame and \(L_{\rm sh}\) for the sharp point lattice. For the last-line generator \(t_z\) and the central weight differences \(s_\tau=|\lambda_\tau|-|\sigma_\tau|\), the normalized restriction satisfies
\[
 {\rm br}_{t_z}=\left(\prod_\tau\tau(t_z)^{-s_\tau}\right){\rm br}_{\rm std},\qquad {\rm br}_{t_z}(L_{\rm st})\subseteq L_{\rm sh}.
\]
 It is covariant in the terminal class and commutes
with the reference and twisted evaluations at the respective tuples.
\end{lemma}

\begin{proof}
Use Proposition~\ref{prop:integral-terminal-branching} in the
standard coordinates before transporting the class bases.
At a terminal class let \(L'\) be its sharp lattice. In the
lifted class the decorated lattice is
\[
  \widetilde L=\{(v',v_N): v'\in L',\quad t_zv_N\in O_E\}.
\]
Here \(t_z\) is the last-coordinate multiplier of the
representative. Fix a rational prime \(p\), and choose a basis
matrix \(B'\) of \(L'\) at the places over \(p\). A single
matrix over \(E\) suffices, since the localization of the
Dedekind domain at this finite set is semilocal and its
fractional ideals are principal. A basis matrix of
\(\widetilde L\) there is \(B=\operatorname{diag}(B',t_z^{-1})\).

Let \(s_\tau=\sum_i\lambda_{\tau i}-\sum_i\sigma_{\tau i}\)
using the full weights. The actual restriction at the
representative uses
\[
   \mathrm{br}_{t_z}=
        \Big(\prod_\tau\tau(t_z)^{-s_\tau}\Big)\mathrm{br}_{\rm std}.
\]
Indeed the last diagonal group has character \(a^{s_\tau}\)
on the selected summand by the two central weights. In the
inverse-transpose representation \(\rho^\vee\), reading in
the inverse basis uses \(B^t\) on the polynomial variables.
The block identity is therefore
\begin{equation}\label{eq:terminal-block-identity}
 \rho^\vee_\sigma((B')^{-1})\,\mathrm{br}_{t_z}(G)
       =\mathrm{br}_{\rm std}\big(\rho^\vee_\lambda(B^{-1})G\big).
\end{equation}
In particular the last diagonal entry acts with the full
weight difference, even when the polynomial degree has
been reduced by a determinant twist.

At the lifted class evaluate an integral stage form. A change
from its chosen representative acts on its value by the same
lattice transport as on its decorated lattice, times the finite
character, a unit. Thus its coordinates in \(B\) belong to
the local Weyl lattice. Write \(F(z)\) for the value at the lifted class and
\(T(z')\) for its restriction at the matched terminal class.
By the stated restriction on point values,
\[
   \rho^\vee_\sigma((B')^{-1})T(z')
      =\mathrm{br}_{\rm std}\bigl(\rho^\vee_\lambda(B^{-1})F(z)\bigr).
\]
The right side lies in the standard integral Weyl lattice
by integral lowering. The finite character of the restriction
is the quotient of the two class characters. The same block
intertwining identity with the stabilizer action gives covariance
under a change of representative. Equivalently this is the
restriction of the complex automorphic form after embedding
the coefficients, and injectivity of the embedding gives the
covariance over \(K\).

A covariant family on a finite set of class representatives
with values in their \(K\)-spaces defines the rational terminal
point. Its localization is in the integral point lattice
precisely when all these values are locally integral.
For finitely many local values \(y_c\) with clearing factors \(s_c\) away from \(w\), a single factor works, since
\[
  s=\prod_c s_c,\qquad v_w(s)=0,\qquad sy_c=(s/s_c)(s_c y_c).
\]
 To
see the reverse inclusion, clear their finitely many
denominators by a single element outside the place. This
preserves the linear covariance equations. Intersecting over
the places gives the assertion for integral global inputs.
These comparisons use the lifted classes and \(t_z\) of the
tuple in question. They commute with reference or twisted
evaluation on that tuple. \label{cor:terminal-realization-3}
\end{proof}

\paragraph{Realizing a sharp test.}
The selected row uses the coordinate image of the integral
level lattice. Here is the descent to evaluations in that
lattice. Let \(j\) be the semilinear embedding of level
coordinates in the sharp complex row ambient space at this
tuple. Restricting an algebraic level form to class
representatives identifies it with its coordinates satisfying
the stabilizer equations. Conversely, if \(a\) is a chosen representative with
assigned value \(y_a\), the covariance operator \(R(g)\)
assigns \(R(g)y_a\) to its translate.

When \(ga=g'a\),
the stabilizer equation \(R((g')^{-1}g)y_a=y_a\)
makes these assignments agree.
Indeed covariance gives
\[
  R(g)y_a=R(g')R((g')^{-1}g)y_a=R(g')y_a .
\]
 This produces exactly the
algebraic form with those coordinates. The integral coordinate lattice
consists of the vectors whose values lie in the transported
flag spans at the level points. This description also results
if one first uses the full finite level action and then
restricts to the given character: a vector already in the
character space is recovered from its level coordinates by
inclusion.
In particular, for a character \(\chi\) of the level group, the integral points of its character space are given in level coordinates by
\[
  M[\chi]=\{v\in M:R(h)v=\chi(h)v\ \text{for every level element }h\}.
\]
 No projector denominator enters this operation.

Choose the designated finite generators \(v\) of the sharp
ambient lattice. In the selected row data they are the
embedded coordinate generators of the integral level lattice.
More explicitly, each has a coordinate lift \(g_v\) with
\(j(g_v)=v\), and a level coordinate \(y\) is integral exactly
when \(j(y)\) is an \(\mathcal O_K\)-combination of these
generators under the coefficient embedding. This description
follows by taking a finite generating set of the integral
coordinate module before extension of coefficients, and
transporting its span with the chosen level isomorphism.
If \(F=\sum_v\iota(c_v)v\) is a sharp integral test, take

\[
  y=\sum_v c_vg_v,
\]
 with the coefficients mapped to
\(\mathcal O_K\) if the record uses its original coefficient ring.
By semilinearity,
\[
          j(y)=\sum_v \iota(c_v)j(g_v)
                     =\sum_v\iota(c_v)v=F.
\]
Each generator \(g_v\) is in the integral level module and
so is \(y\). The stabilizer equations are linear, so the
covariant extension from \(y\) is integral at every class.
This uses the image of the integral module, including its
index in the rational coordinates.

Read that form at the terminal representative using
Lemma~\ref{lem:terminal-integral-class-frame}. By the decorated
lattice identification its value is in the flag span for
\(L'\). The sharp ambient identification and the inverse
period identification give the same function there: both
extend the same level coordinates by covariance.
Indeed if two covariant extensions \(F_1,F_2\) agree at the representative \(a\), then on its entire class
\[
  F_1(ga)=R(g)F_1(a)=R(g)F_2(a)=F_2(ga).
\]
 This applies
to the reference identification and to the twisted one with
its own representative. Thus \emph{every} sharp test is
represented at every terminal class by an integral row
polynomial for the lattice dual to the column just
constructed. The pairing in the terminal sum can now be
applied class by class. \label{cor:terminal-realization-4}
